% !TeX root = main.tex
% !TeX program = xelatex
% Include this chapter after draft_zh/02_授权流程与分析目标.
% Chinese compilation uses the existing IEICE class with XeLaTeX and xeCJK.
% Results below are tied to archived evidence. No new verification is claimed.

\section{模型与分析结果}
\label{sec:models-analysis-results}

本章沿第二章的授权流程，分析策略求值、OTK 发放与使用以及 ACT 并发准入三类约束。
第~\ref{subsec:policy-evaluation-consistency}~节对照参考语义与实际匹配函数，并考察决策分歧对凭证签发的影响。
第~\ref{subsec:otk-issuance-use}~节采用 Tamarin 检查发放与使用事件的身份对应，
第~\ref{subsec:act-concurrent-quota}~节采用 TLA$^{+}$/TLC 分析共享额度的交错执行，并与受控原生方法实验对照。
实现分析依据固定版本的参考代码~\cite{sagaReferenceImpl2026}，
两项形式化模型均由相关代码手工抽象，保留与相应性质有关的身份、状态及执行步骤。
表~\ref{tab:authorization-analysis-results}~概括各项分析的依据与结果。

\begin{table}[tb]
  \caption{三类授权约束的分析依据与结果}
  \label{tab:authorization-analysis-results}
  \centering
  \small
  \begin{tabular}{@{}p{0.18\linewidth}p{0.26\linewidth}p{0.44\linewidth}@{}}
    \hline
    检查目标 & 分析依据 & 主要结果 \\
    \hline
    策略求值 & 函数对照与受控执行 & 规则顺序改变预算，并传播至 ACT 签发与存储 \\
    OTK 发放与使用 & Tamarin\newline 固定配置模型 & 发放上限与单次消费成立，同人使用对应出现反例 \\
    ACT 累计准入 & TLC 有限状态检查与原生方法实验 & 双连接共享末次额度时出现超额准入 \\
    \hline
  \end{tabular}
\end{table}

\subsection{策略求值一致性}
\label{subsec:policy-evaluation-consistency}

策略分析检验规则优先语义能否准确落实为实际预算决定。
对接触策略 $CP_A$ 和发送方身份 $s$，记实际函数与参考求值分别为
$\operatorname{Eval}_{\mathrm{impl}}$ 和 $\operatorname{Eval}_{\mathrm{ref}}$，检查目标为
\begin{equation}
  \operatorname{Eval}_{\mathrm{impl}}(CP_A,s)
  = \operatorname{Eval}_{\mathrm{ref}}(CP_A,s).
  \label{eq:policy-evaluation-consistency}
\end{equation}
参考求值依据协议的最具体匹配规则要求及匹配器的具体性评分确定~\cite{syros2025sagaV2,sagaReferenceImpl2026}。
本节采用合法且具有唯一最高评分匹配规则的配置。

% Evidence: proofs/evidence/authz-stage1-20260916T123834944479Z/policy-matcher-sanity.json
规则顺序对照给出了式~\eqref{eq:policy-evaluation-consistency}~的具体反例。
对于身份 \texttt{alice@\allowbreak example.com:\allowbreak agent}，规则
\texttt{alice@\allowbreak example.com:*} 的预算为 $-1$，通配规则 \texttt{*} 的预算为 $10$，
两者的具体性评分分别为 $70$ 和 $0$。
参考结果均为 $-1$。
直接调用原始匹配函数时，具体拒绝规则在前、通配允许规则在后得到 $10$，交换顺序则得到 $-1$。
该差异源于 \texttt{best\_pattern} 的状态维护。
该变量初始化为 \texttt{None}，对应评分为 $-1$。
循环更新预算时未更新该变量，使后续匹配持续与初始评分比较，通配允许规则因而能够覆盖先前的具体拒绝规则。

% Evidence: proofs/consumer_validation/evidence/20260926T070021921746Z-d64c8064/cases/da/
% Scope: production access route and receiver handler, ending at ACT issuance/storage.
受控消费者执行进一步显示，规则选择分歧可以传播至后续 ACT 签发。
在具体拒绝规则先于通配允许规则的配置中，Provider 与接收方的匹配结果均为 $10$，
执行完成 OTK 发放与消费，并发送及保存 ACT。
交换规则顺序后，Provider 返回拒绝，执行在材料发放之前终止。
实验保留生产授权函数、数据库更新、签名校验及令牌密码操作，
以本地身份材料与替代传输组织执行，并在会话处理入口结束观察。
该结果将策略求值分歧的影响定位到凭证签发，也说明两端复用同一匹配实现时，两次检查会继承同一规则选择偏差。

\subsection{OTK 发放与使用}
\label{subsec:otk-issuance-use}

OTK 分析分别检验原有发放上限与额外的同人使用对应关系。
对获准发送方 $s$，记任一执行前缀中 $P$ 向其发放 $A$ 的 OTK 的次数为 $I_{A,s}$，初始预算为 $\beta_{A,s}$。
原有发放约束为
\begin{equation}
  I_{A,s} \leq \beta_{A,s}.
  \label{eq:otk-issuance-bound}
\end{equation}
进一步记 $\mathsf{Issue}(A,s,k)$ 为 $P$ 向 $s$ 发放 OTK $k$ 的事件，
$\mathsf{Use}(A,s,k)$ 为 $s$ 取得并解密该 OTK 对应的新 ACT 后，其首个应用请求获 $A$ 准入的事件。
对任意 $s,k$ 及事件位置 $i$，本文额外检查
\begin{equation}
  \begin{aligned}
    &\mathsf{Use}(A,s,k)@i \\
    &\quad\Rightarrow\exists j<i.\ \mathsf{Issue}(A,s,k)@j.
  \end{aligned}
  \label{eq:otk-same-sender-correspondence}
\end{equation}
式~\eqref{eq:otk-same-sender-correspondence}~表达从领取到成功使用的同人对应，
用于检验将领取预算进一步解释为按身份限制新许可次数时的约束。

% Model/source map: proofs/otk_transfer/saga_otk_transfer.spthy and SOURCE_MAP.zh-CN.md
% Use is the RequestAccepted(A,s,k,t,pac) projection, after issuance and decryption.
Tamarin 模型保留了发放计数、两处 OTK 状态与实际连接者的身份及密钥依赖~\cite{meier2013tamarin}。
模型固定 $P$、接收方 $A$ 和两个发送方 $B,C$，设置两枚 OTK，且 $\beta_{A,B}=\beta_{A,C}=1$。
两发送方均获接触策略允许，并各自具有正确绑定的身份和访问控制公钥。
Provider 的剩余领取额度、尚未发放 OTK 与 $A$ 的尚未消费 OTK 分别表示为独立状态。
模型采用私有端点消息表示 TLS 通信，并保留 DH、密钥派生、ACT 加密与解密操作的符号依赖。
转交行为允许 $B$ 将其已合法领取、尚未消费的 OTK 公钥主动发送给 $C$，长期私钥保持保密。
每枚新 ACT 的使用观察到其首次请求准入。

% Evidence: proofs/otk_transfer/evidence/run-20261025T033815645931Z/forwarding/witness.zh-CN.md
转交见证显示，发放上限与更强的同人对应具有不同的检查结果。
$C$ 先使用自行领取的 $k_C$，取得并解密新 ACT，完成首次请求准入。
其领取额度耗尽后，$P$ 拒绝再次发放。
$B$ 随后将另一枚已领取而未消费的 $k_B$ 转交给 $C$。
$C$ 以自己的身份、访问控制公钥和私钥使用 $k_B$，取得并解密另一枚 ACT，其请求再次获准入。
因此，该执行中 $C$ 的领取次数为 $1$，成功使用的不同 OTK 数量为 $2$。

% Evidence: proofs/otk_transfer/evidence/run-20261025T033815645931Z/summary.json
模型检查在上述固定配置中验证了式~\eqref{eq:otk-issuance-bound}~及每枚 OTK 的单次消费，
同时给出式~\eqref{eq:otk-same-sender-correspondence}~的反例。
仅移除主动转交后，$C$ 的正常流程仍可达，同人对应性质得到验证。
这组对照将差异定位到领取材料的转交与实际使用身份之间的关系。
Provider 的计数限制按领取身份执行，而接收方依据未消费 OTK 与实际连接者的公钥建立新 ACT，
因而将发放预算解释为实际使用上限时，还需明确相应的身份约束。

\subsection{ACT 并发额度执行}
\label{subsec:act-concurrent-quota}

ACT 分析检查共享额度的并发执行能否维持式~\eqref{eq:act-admission-bound}~的累计准入上限。
参考实现的 \texttt{token\_is\_valid()} 在锁内检查活动记录、有效期、剩余额度及发送方公钥，成功返回时释放锁。
\texttt{receive\_conversation()} 随后重新取得同一把锁，将剩余额度 $r$ 更新为 $\max(0,r-1)$，再进入请求处理。
额度检查与扣减位于两个独立的临界区，而不同连接可并发访问同一 ACT 的状态~\cite{sagaReferenceImpl2026}。

% Model/source map: proofs/quota_validation/SagaQuotaDelivery.tla and DELIVERY_MODEL.md
TLA$^{+}$ 模型通过分离检查、扣减与后端调用步骤表达上述执行结构~\cite{lamport2002specifying}。
模型保留双端余额、各连接的局部进度、消息交付状态及独立的准入与后端调用计数。
检查动作读取接收方余额并推进连接位置，扣减动作更新余额并增加准入计数，后端调用由后续动作记录。
本节采用初始额度 $Q(t)=1$、两个连接各至多提交一个请求、非配额元数据有效且故障数为零的有限配置。

% Evidence: proofs/quota_validation/evidence/delivery-20261024T150331844425Z/concurrent_race/
TLC 在该配置中给出了累计准入上限的反例。
当接收方余额为 $1$ 时，两个连接先后通过检查，随后分别执行扣减。
以接收方余额 $r$ 与累计准入次数 $n=N_{\mathrm{adm}}(t)$ 投影关键交错，可表示为
\begin{equation}
  \begin{aligned}
    (1,0)&\xrightarrow{\mathsf{check}_{c_1},\,\mathsf{check}_{c_2}}(1,0)\\
         &\xrightarrow{\mathsf{debit}_{c_1}}(0,1)
          \xrightarrow{\mathsf{debit}_{c_2}}(0,2).
  \end{aligned}
  \label{eq:act-last-credit-interleaving}
\end{equation}
截断扣减使最终余额保持为 $0$，累计准入次数却达到 $2$。
该坏前缀的后端调用计数仍为 $0$，反例定位于准入阶段。

% Evidence: output/analysis/qmax-20261020/results.json and probe_qmax.py
受控原生方法实验独立观察到了同类交错引起的额外后端调用。
初始额度为 $50$ 时，串行提交 $51$ 条请求产生 $50$ 次计数后端调用，余额为 $0$。
先串行处理 $49$ 条，再以线程屏障使两个连接均通过剩余一次额度的原生检查，
随后释放线程继续执行，可得到累计 $51$ 次后端调用，余额仍为 $0$。
三次受控重复记录了相同结果。
实验保留原生身份校验、令牌处理与额度扣减，使用本地证书发现、替代传输和计数后端组织执行。
有限状态反例定位了准入上限的失效步骤，方法实验则在原生执行中观察了相同检查与扣减交错的后续调用结果。
